% Options for packages loaded elsewhere
\PassOptionsToPackage{unicode}{hyperref}
\PassOptionsToPackage{hyphens}{url}
\documentclass[
  ignorenonframetext,
]{beamer}
\newif\ifbibliography
\usepackage{pgfpages}
\setbeamertemplate{caption}[numbered]
\setbeamertemplate{caption label separator}{: }
\setbeamercolor{caption name}{fg=normal text.fg}
\beamertemplatenavigationsymbolsempty
% remove section numbering
\setbeamertemplate{part page}{
  \centering
  \begin{beamercolorbox}[sep=16pt,center]{part title}
    \usebeamerfont{part title}\insertpart\par
  \end{beamercolorbox}
}
\setbeamertemplate{section page}{
  \centering
  \begin{beamercolorbox}[sep=12pt,center]{section title}
    \usebeamerfont{section title}\insertsection\par
  \end{beamercolorbox}
}
\setbeamertemplate{subsection page}{
  \centering
  \begin{beamercolorbox}[sep=8pt,center]{subsection title}
    \usebeamerfont{subsection title}\insertsubsection\par
  \end{beamercolorbox}
}
% Prevent slide breaks in the middle of a paragraph
\widowpenalties 1 10000
\raggedbottom
\AtBeginPart{
  \frame{\partpage}
}
\AtBeginSection{
  \ifbibliography
  \else
    \frame{\sectionpage}
  \fi
}
\AtBeginSubsection{
  \frame{\subsectionpage}
}
\usepackage{iftex}
\ifPDFTeX
  \usepackage[T1]{fontenc}
  \usepackage[utf8]{inputenc}
  \usepackage{textcomp} % provide euro and other symbols
\else % if luatex or xetex
  \usepackage{unicode-math} % this also loads fontspec
  \defaultfontfeatures{Scale=MatchLowercase}
  \defaultfontfeatures[\rmfamily]{Ligatures=TeX,Scale=1}
\fi
\usepackage{lmodern}
\ifPDFTeX\else
  % xetex/luatex font selection
\fi
% Use upquote if available, for straight quotes in verbatim environments
\IfFileExists{upquote.sty}{\usepackage{upquote}}{}
\IfFileExists{microtype.sty}{% use microtype if available
  \usepackage[]{microtype}
  \UseMicrotypeSet[protrusion]{basicmath} % disable protrusion for tt fonts
}{}
\makeatletter
\@ifundefined{KOMAClassName}{% if non-KOMA class
  \IfFileExists{parskip.sty}{%
    \usepackage{parskip}
  }{% else
    \setlength{\parindent}{0pt}
    \setlength{\parskip}{6pt plus 2pt minus 1pt}}
}{% if KOMA class
  \KOMAoptions{parskip=half}}
\makeatother
\setlength{\emergencystretch}{3em} % prevent overfull lines
\providecommand{\tightlist}{%
  \setlength{\itemsep}{0pt}\setlength{\parskip}{0pt}}
\usepackage{bookmark}
\IfFileExists{xurl.sty}{\usepackage{xurl}}{} % add URL line breaks if available
\urlstyle{same}
\hypersetup{
  hidelinks,
  pdfcreator={LaTeX via pandoc}}

\author{\texorpdfstring{}{}}
\date{}

\begin{document}

\begin{frame}[fragile]{Supplement: Limitations and assumptions}
\protect\phantomsection\label{supplement-limitations-and-assumptions}
This exemplar prioritizes \textbf{reproducibility and teachability} over
geographic realism. Riverbend is fictional; weights and source labels
illustrate schema use, not empirical findings about any jurisdiction.

\textbf{No live data.} The pipeline does not call HTTP APIs, scrape
documents, or perform NLP on attachments. Every number in the handbook
narrative traces to a hand-authored YAML row or to derived metrics
computed from those rows.

\textbf{Weights are subjective.} Reviewers assign \texttt{weight} in
\([0,1]\) without a universal rubric. Cross-team calibration is a
governance task outside this repository; the code assumes weights are
already consensus values.

\textbf{Single corpus per build.} The scripts load one fixture path
(\texttt{data/fixtures/riverbend\_area.yaml} in production use of this
demo). Merging multiple corpora or federating regions would require new
modules and tests, not configuration toggles alone.

\textbf{Outline static in code.} \texttt{HANDBOOK\_TEMPLATE} is Python
source, not a per-area YAML file. Customizing chapters without editing
code is possible only within the existing section files and themes;
adding a ninth scored section demands code and test updates.

\textbf{Citation depth.} \texttt{source\_label} is a short string, not a
BibTeX key. Formal citations belong in manuscript markdown using
\texttt{references.bib}; the corpus stays lean {[}@bowker2005{]}.

\textbf{Threshold semantics.} ``Gap'' means strictly below threshold,
not ``close.'' Sections barely above the line still pass; policy teams
may want stricter thresholds or secondary rules implemented in future
metrics fields.
\end{frame}

\end{document}
